\section{Local and rational isometries of tensor squares}
\label{local:section}

Let $q_1,q_2$ be the polynomials of Corollary~\ref{seed:prepared},
in the same even number $s$ of variables.  Write $L_j$ for their Milnor
lattices and $S_j$ for their integral Seifert forms.  The two inputs from
that corollary are an isomorphism of the integral monodromy modules
and an isomorphism of the real Seifert forms.  We now combine these
inputs, first locally and then rationally.

Choose an odd integer $m$ divisible by the orders of both monodromies,
and choose $r\in\Z$ with $2r\equiv1\pmod m$.  With $s$ even and the
monodromy convention in Equation~\eqref{eq:bg-monodromy}, put
\begin{equation}\label{local:symmetric-data}
 T_j=S_j^{-1}S_j^{\mathsf t},\qquad B_j=S_jT_j^r.
\end{equation}
Throughout, an \emph{equivariant isometry} between data $(L,T,B)$
means an isometry of the forms that intertwines the displayed operators.
Tensor products carry both the tensor product form and the tensor product
operator.

\begin{lemma}[Symmetrization]\label{local:symmetrize}
Each $B_j$ is an integral unimodular symmetric form preserved by $T_j$.
The pair $(B_j,T_j)$ recovers $S_j$ by $S_j=B_jT_j^{-r}$.
The real Seifert isometry is an equivariant real isometry of the
symmetric data.  Symmetrization commutes with tensor products of
even-variable factors, up to the common signs in
Theorem~\ref{bg:ts}.
\end{lemma}

\begin{proof}
Suppress the subscript.  The defining equality $S^{\mathsf t}=ST$ implies
\[
 T^{\mathsf t}ST=S,\qquad
 B^{\mathsf t}=(T^r)^{\mathsf t}S^{\mathsf t}
             =ST^{1-r}=ST^r=B.
\]
It also gives $T^{\mathsf t}BT=B$.  Integrality and unimodularity follow
from those of $S$ and $T$.  An isometry of $S$ intertwines
$S^{-1}S^{\mathsf t}$, and therefore preserves $B$ as well.
For tensor products, use the same $r$ modulo a common odd exponent and
the identity
\[
 (S_1\otimes S_2)(T_1\otimes T_2)^r
       =(S_1T_1^r)\otimes(S_2T_2^r).
\]
The overall Seifert signs supplied by the integral Thom--Sebastiani
theorem depend only on the factor dimensions~\cite{Sakamoto1974};
they are thus the same on the two sides of every comparison here.
\end{proof}

The local comparison will concern the first tensor squares
\[
 A_j=(L_j,T_j,B_j)^{\otimes2}.
\]
Our goal is an equivariant isometry between $A_1$ and $A_2$ over every
$\Z_\ell$, together with one over $\Q$.  The integral operator-module
isomorphism lets us compare the two forms on a common lattice.  At odd
primes a tensor square makes these forms isometric; at $2$ the prepared
operators already give uniqueness before tensoring.  The real Seifert
isometry then supplies the archimedean input for rational descent.

\subsection{Odd primes}

We first give the elementary lifting fact needed for a simultaneous
isometry of the form and operator.

\begin{lemma}[Lifting in a finite operator algebra]\label{local:lifting}
Let $\ell$ be an odd prime and let $\mathcal A$ be a commutative
$\Z_\ell$-algebra that is finite as a $\Z_\ell$-module.  Write $J$ for
its Jacobson radical.  Then $\mathcal A/J$ is a product of finite fields,
and $\mathcal A$ is complete for the $J$-adic topology.  For a unit
$u\in\mathcal A$:
\begin{enumerate}
\item if its image in every residue field is a square, then $u$ is a
square in $\mathcal A$;
\item $u=a^2+b^2$ for some $a,b\in\mathcal A$.
\end{enumerate}
\end{lemma}

\begin{proof}
The algebra is complete for the $\ell$-adic topology.  For every
$c\in\mathcal A$, the geometric series for $(1-\ell c)^{-1}$ converges,
so $\ell\mathcal A\subseteq J$.  The finite algebra
$\mathcal A/\ell\mathcal A$ has nilpotent radical and semisimple quotient
a product of finite fields.  Consequently $J^N\subseteq\ell\mathcal A$
for some $N$, proving that the two topologies agree and establishing the
first assertions.

For (1), choose a lift $v$ of a square root in $\mathcal A/J$.
Both $v$ and $2v$ are units.  If $e=v^2-u\in J^k$, replace $v$ by
$v-e/(2v)$.  The new error is $e^2/(4v^2)\in J^{2k}$.
Iteration converges to a square root of $u$.

In a finite field $k$ of odd cardinality, the sets of squares and of
$\bar u$ minus the squares each have $(|k|+1)/2$ elements.  They
intersect, so $\bar u=\bar a^2+\bar b^2$.  For (2), make such a
choice in each residue field and lift it to $a_0,b_0\in\mathcal A$.
Then $c=a_0^2+b_0^2$ is a unit and $u/c\in1+J$.  By (1), write
$u/c=v^2$.  Taking $a=va_0$ and $b=vb_0$ proves the claim.
\end{proof}

\begin{proposition}[Odd-prime tensor squares]\label{local:odd-primes}
Let $L$ be a finite free $\Z_\ell$-module, with $\ell$ odd, and let
$T$ be an automorphism of $L$.  If $B_1,B_2$ are unimodular symmetric
forms preserved by $T$, then
\[
 (L\otimes L,T\otimes T,B_1\otimes B_1)
 \cong
 (L\otimes L,T\otimes T,B_2\otimes B_2)
\]
by an equivariant integral isometry.
\end{proposition}

\begin{proof}
Set $P=B_1^{-1}B_2$, so
$B_2(x,y)=B_1(x,Py)$.  Unimodularity makes $P$ an integral
automorphism.  Symmetry of $B_2$ makes it self-adjoint for $B_1$, and
invariance of both forms gives $PT=TP$.

Write the characteristic polynomial of $P$ modulo $\ell$ as
$\prod_i f_i^{e_i}$, with distinct monic irreducible $f_i$.
Lift this coprime factorization to
$\det(X-P)=\prod_i F_i$ over $\Z_\ell$, where
$F_i\bmod\ell=f_i^{e_i}$.  For completeness, the lift is obtained
successively modulo $\ell^n$.  For two coprime monic residue factors
$f,g$ of degrees $a,b$, the correction map
\[
 (u,v)\longmapsto ug+fv,\qquad \deg u<a,\quad\deg v<b,
\]
is an isomorphism onto the polynomials of degree less than $a+b$:
its kernel is zero by coprimality and the dimensions agree.  It
therefore corrects the product by any prescribed error at the next
power of $\ell$, without changing monicity.  Completeness gives the
lift, and iteration handles any number of factors.  The Chinese remainder
idempotents for the $F_i$, evaluated at $P$, give
\[
 L=\bigoplus_i L_i,\qquad F_i(P_i)=0,
 \qquad P_i=P|_{L_i}.
\]
These idempotents are polynomials in $P$.  They are self-adjoint and
commute with $T$, so the sum is orthogonal for $B_1$ and $B_2$ and each
$L_i$ is $T$-stable.  Each restricted form is unimodular.

On $L_i\otimes L_j$, let $\mathcal A_{ij}$ be the algebra generated
by $P_i\otimes1$ and $1\otimes P_j$.  Cayley--Hamilton bounds the
degrees in these two generators, so this is a finite commutative
$\Z_\ell$-algebra of endomorphisms.  Every element is self-adjoint
for the restricted form $B_1\otimes B_1$ and commutes with
$T\otimes T$.  The unit
\[
 U_{ij}=P_i\otimes P_j
\]
relates the two restricted tensor forms.  Its inverse belongs to
$\mathcal A_{ij}$, since the inverses of the generators are polynomials
in them with integral coefficients.

If $i=j$, consider a residue field of $\mathcal A_{ii}$.
The relations $F_i(P_i)=0$ imply that the two generator images
$x,y$ are roots of $f_i$.  They are nonzero because the generators
are units.  In a finite field, two roots of the same irreducible
polynomial over $\F_\ell$ are Frobenius conjugates.  Hence
$y=x^{\ell^h}$ for some $h\geq0$, and
\[
 xy=x^{1+\ell^h}
    =\left(x^{(1+\ell^h)/2}\right)^2.
\]
Lemma~\ref{local:lifting} supplies $a\in\mathcal A_{ii}$ with
$a^2=U_{ii}$.  Because $a$ is self-adjoint,
\[
 (B_1\otimes B_1)(ax,ay)=(B_2\otimes B_2)(x,y)
 \qquad(x,y\in L_i\otimes L_i).
\]

For $i<j$, pair $L_i\otimes L_j$ with $L_j\otimes L_i$.
The factor swap identifies their operators and both their forms.
On the resulting two equal blocks, use
Lemma~\ref{local:lifting} to choose $a,b\in\mathcal A_{ij}$ with
$a^2+b^2=U_{ij}$ and set
\[
 M=\begin{pmatrix}a&b\\-b&a\end{pmatrix}.
\]
Writing $*$ for adjoint with respect to the first tensor form on
the two blocks gives
\[
 M^*M=\begin{pmatrix}U_{ij}&0\\0&U_{ij}\end{pmatrix}.
\]
This is invertible, and $M$ commutes with the tensor operator.
The diagonal-block maps and these paired-block maps combine to an
integral equivariant isometry from the second tensor form to the first.
Its inverse has the direction asserted in the proposition.
\end{proof}

\subsection{The prime two}

Here the odd order and absence of eigenvalue $1$ give a stronger result.
The unramified hermitian classification underlying this step belongs
to the classical local theory developed by Jacobowitz~\cite{Jacobowitz1962}.
We give the needed argument directly, including residue characteristic two.

\begin{proposition}[Dyadic uniqueness]\label{local:dyadic}
Let $L$ be a finite free $\Z_2$-module, and let $T$ satisfy $T^m=1$
for an odd $m$.  Suppose $1$ is not an eigenvalue of $T$ on
$L\otimes\Q_2$.  Any two $T$-invariant unimodular symmetric forms
on $L$ are equivariantly isometric over $\Z_2$.
\end{proposition}

\begin{proof}
The reduction of $X^m-1$ modulo $2$ is separable, since its derivative
is $X^{m-1}$ there.  Lifting its distinct irreducible factors gives
\[
 \mathcal O=\Z_2[X]/(X^m-1)=\prod_i\mathcal O_i,
\]
where each $\mathcal O_i$ is the integer ring of an unramified
extension of $\Q_2$.  The action $X\mapsto T$ decomposes
$L=\bigoplus_i M_i$.  Every $M_i$ is free over $\mathcal O_i$:
it is finite and torsion-free over this discrete valuation ring.

Invariance of a form $B$ says that the adjoint of $X$ is $X^{-1}$.
Let $\sigma$ denote this involution of $\mathcal O$.  If it exchanges
the factors indexed by $i$ and $i'$, then $M_i$ and $M_{i'}$ are
isotropic and pair perfectly with each other; they are orthogonal to
all other factors.  Such a pairing has a unique equivariant isometry
class for the prescribed modules.  Indeed, it identifies $M_{i'}$
with the $\Z_2$-dual of $M_i$, with its $\mathcal O_i$-action
transported by $\sigma$.  An isomorphism on $M_i$ extends uniquely by
its inverse dual to an isometry of the pair.

It remains to treat a nonzero $M_i$ whose factor is preserved by
$\sigma$.  Write $\mathcal O_i=\mathcal D$ and
$\mathcal D_0=\mathcal D^\sigma$.
The involution on $\mathcal D$ is nontrivial.  Otherwise the image
$\zeta$ of $X$ would satisfy $\zeta=\zeta^{-1}$; its odd order would
force $\zeta=1$, contrary to the hypothesis.  Thus
$\mathcal D/\mathcal D_0$ is an unramified quadratic extension.

We record explicitly how $B$ becomes an integral hermitian form.
The trace pairing of $\mathcal D$ over $\Z_2$ is perfect: modulo
$2$ it is the nondegenerate trace pairing of a finite separable field
extension, and its determinant is therefore a unit.  Hence there is
a unique $\mathcal D$-valued form $h$ such that
\begin{equation}\label{local:trace}
 \operatorname{Tr}_{\mathcal D/\Z_2}\bigl(a h(x,y)\bigr)
       =B(ax,y)\qquad(a\in\mathcal D).
\end{equation}
The adjoint relation and symmetry of $B$ give
\[
 h(ax,by)=a\sigma(b)h(x,y),\qquad
 h(y,x)=\sigma(h(x,y)).
\]
Thus $h$ is hermitian, with linear first variable.
It is unimodular: after extending to fraction fields, the $B$-dual
of $M_i$ equals its $h$-dual.  Indeed, if $B(x,M_i)\subseteq\Z_2$,
then for every $a\in\mathcal D$ and $y\in M_i$,
\[
 B(ax,y)=B(x,\sigma(a)y)\in\Z_2.
\]
Testing all $a$ in~\eqref{local:trace} and using the perfect trace
pairing gives $h(x,M_i)\subseteq\mathcal D$; the converse is immediate.
This argument takes no quotient by $2$.

Every unimodular hermitian form over
$\mathcal D/\mathcal D_0$ is determined by its rank.  Here are the
details in residue characteristic two.  Let $k/k_0$ be the quadratic
residue extension and let $\bar h$ be the reduced hermitian form.
There is a vector of nonzero length.  If all lengths vanished, then
for all $x,y$ and $a\in k$,
\[
 0=\bar h(x+ay,x+ay)
   =\operatorname{Tr}_{k/k_0}
       \bigl(\sigma(a)\bar h(x,y)\bigr).
\]
The finite-field trace pairing would force $\bar h=0$, contradicting
nonsingularity.  A lift of a vector of nonzero length has unit length
$u\in\mathcal D_0^\times$.  Its line splits off integrally, by the
projection $x\mapsto u^{-1}h(x,v)v$.  Its orthogonal complement is
again unimodular.  Induction diagonalizes $h$ with all diagonal
entries in $\mathcal D_0^\times$.

Every such unit is a norm from $\mathcal D^\times$.  The residue norm
$k^\times\to k_0^\times$ is surjective.  After lifting a residue
preimage, improve a norm congruence modulo $2^n$ to one modulo
$2^{n+1}$ using
\[
 N_{\mathcal D/\mathcal D_0}(1+2^nz)
  \equiv1+2^n\operatorname{Tr}_{\mathcal D/\mathcal D_0}(z)
       \pmod{2^{n+1}},\qquad n\geq1.
\]
The trace modulo $2$ is surjective, so the needed $z$ always exists.
The corrections converge.  Rescaling each diagonal vector therefore
makes every diagonal entry $1$.  The ranks of the $M_i$ depend only
on the operator module, so this proves uniqueness on every fixed
factor and completes the proof.
\end{proof}

\subsection{Rational descent}

Although neither $T_j$ has eigenvalue $1$, its tensor square can:
an eigenvalue and its inverse have product $1$.  The rational argument
must therefore include an ordinary quadratic-form component at
eigenvalue $1$, as well as the hermitian components for nonreal
eigenvalues.  The odd order still excludes $-1$.

We shall use Hasse--Minkowski in the form stated in
Theorem~\ref{bg:hasse}.  The passage from quadratic forms to hermitian
forms needed here has an elementary proof.

\begin{lemma}[Hermitian forms from their lengths]
\label{local:hermitian-lengths}
Let $E/F$ be a quadratic field extension with $\operatorname{char}F\ne2$,
and let $h_1,h_2$ be nonsingular hermitian forms.  If the quadratic
forms
\[
 Q_{h_j}(x)=h_j(x,x)
\]
on the underlying $F$-spaces are isometric, then $h_1,h_2$ are
$E$-linearly isometric.
\end{lemma}

\begin{proof}
The polar pairing of $Q_h$ is
$\tfrac12\operatorname{Tr}_{E/F}h(x,y)$.
If it annihilates $x$, testing scalar multiples of $y$ and using
the nondegenerate trace pairing shows $h(x,y)=0$ for every $y$.
Thus $Q_h$ is nonsingular.  In particular a nonzero hermitian space
contains a vector of nonzero length, and such a vector splits off
an orthogonal hermitian line.

The quadratic cancellation used below is the reflection argument
underlying Witt cancellation~\cite{Witt1937}; we include the needed
form of it.  For a quadratic
form $Q$ with polar form $b$ and a vector $z$ of nonzero length, the
reflection
\[
 \rho_z(v)=v-\frac{2b(v,z)}{Q(z)}z
\]
is an isometry, as direct expansion shows.  If $Q(x)=Q(e)=a\ne0$,
then $Q(x-e)+Q(x+e)=4a$.  If $Q(x-e)\ne0$, the reflection in $x-e$
sends $x$ to $e$.  Otherwise reflection in $x+e$ sends $x$ to $-e$,
and reflection in $e$ sends $-e$ to $e$.  It follows that an isometry
$\langle a\rangle\perp V\cong\langle a\rangle\perp W$ can be
adjusted to match the two displayed vectors, and then restricts to
an isometry $V\cong W$.

Now write $E=F(\sqrt d)$.  Choose $x$ of hermitian length $a\ne0$
for $h_1$, and let $y$ be its image under an isometry of the length
quadratic forms.  The lines $Ex$ and $Ey$ are hermitian isometric,
and their length quadratic forms are both
\[
 aN_{E/F}=\langle a,-ad\rangle.
\]
Split these hermitian lines off.  Cancel the two nonzero quadratic
lines $\langle a\rangle$ and $\langle-ad\rangle$ successively by
the preceding reflection argument.  The complementary hermitian
forms again have isometric length forms.  Induction on their
$E$-dimension proves the lemma.
\end{proof}

\begin{lemma}[Equivariant rational local--global principle]
\label{local:hasse-descent}
Let $(V_j,U_j,\beta_j)$ be rational vector spaces with nonsingular
symmetric forms preserved by operators of finite odd order.
If they are equivariantly isometric over $\R$ and over $\Q_\ell$
for every prime $\ell$, then they are equivariantly isometric over $\Q$.
\end{lemma}

\begin{proof}
Take a common odd exponent for the operators.  The cyclotomic
decomposition of each $V_j$ is orthogonal: its idempotents are
unchanged under $U_j\mapsto U_j^{-1}$, which is the adjoint
involution.  Each equivariant local isometry respects these
decompositions, and corresponding component dimensions agree.

The eigenvalue-$1$ component carries an ordinary rational symmetric
form.  Its local isometries give a rational isometry by
Theorem~\ref{bg:hasse}.  Odd order excludes eigenvalue $-1$.
For any remaining cyclotomic component put
\[
 E=\Q(\zeta_d),\qquad F=\Q(\zeta_d+\zeta_d^{-1}),
 \qquad d>1.
\]
The component is an $E$-vector space, with $U_j$ acting by $\zeta_d$.
Complex conjugation $\sigma$ is a nontrivial involution of $E$.
The same trace construction as~\eqref{local:trace}, now over $\Q$,
gives a unique nonsingular $E/F$-hermitian form $h_j$ satisfying
\[
 \operatorname{Tr}_{E/\Q}\bigl(a h_j(x,y)\bigr)
      =\beta_j(ax,y)\qquad(a\in E).
\]

An equivariant isometry over a rational completion is linear over
the corresponding algebra $E\otimes_\Q\Q_\ell$, or over
$E\otimes_\Q\R$, and preserves the trace identity.  Nondegeneracy
of the extended trace pairing therefore makes it an isometry of
the extended hermitian forms.  Decomposing
$F\otimes_\Q\Q_\ell$ into its completion factors supplies an
isometry at every place of $F$ above $\ell$; decomposing
$F\otimes_\Q\R$ supplies every archimedean place.  This also
covers finite places where $E/F$ splits, since the trace pairing of
the resulting quadratic \'etale algebra is still nondegenerate.

Consequently the nonsingular quadratic length forms $Q_{h_1}$ and
$Q_{h_2}$ are isometric over every completion of $F$.  Hasse--Minkowski
over $F$ gives a global quadratic isometry
\cite[Theorem~25.64(c)]{Milne2022}.  Lemma~\ref{local:hermitian-lengths} upgrades its
existence to an $E$-linear hermitian isometry.  Such an isometry
preserves the rational symmetric form and intertwines $U_j$.
Combining the cyclotomic components proves the claim.
\end{proof}

\begin{proposition}[Local and rational square isometries]
\label{local:rational}\label{local:conclusion}
Set
\[
 A_j=(L_j,T_j,B_j)^{\otimes2}.
\]
There are an equivariant rational isometry
$R:A_1\otimes\Q\longrightarrow A_2\otimes\Q$ and, for every prime
$\ell$, an equivariant isometry
$D_\ell:A_1\otimes\Z_\ell\longrightarrow A_2\otimes\Z_\ell$.
These can be chosen so that $D_\ell=R$ after extension to $\Q_\ell$
for all but finitely many primes.
\end{proposition}

\begin{proof}
Use the integral operator-module isomorphism of
Corollary~\ref{seed:prepared} to transport the two forms to a common
module with a common operator.  Proposition~\ref{local:odd-primes}
gives the required square isometry at each odd prime.
Proposition~\ref{local:dyadic} applies at $2$, because the prepared
operators have odd order and no eigenvalue $1$; tensor its isometry
with itself.  Lemma~\ref{local:symmetrize} supplies the real isometry,
which also tensors.  Lemma~\ref{local:hasse-descent} now gives $R$.
Finally, the matrices of $R$ and $R^{-1}$ in integral
lattice bases have denominators supported at finitely many primes.
At every other prime $R$ itself is an integral lattice isometry,
so choose $D_\ell=R$ there.
\end{proof}
